\chapter{Kimlik Guvenligi}

\section*{Giris}
Kimlik guvenligi, modern kurumlarda ag sinirlarindan daha kritik bir kontrol katmanidir. Dogru kimlik dogrulama ve yetkilendirme olmadan sifreleme, ag guvenligi ve uygulama guvenligi tam etki uretemez.

\section{Kimlik ve Erisim Yonetimi (IAM) Mimarisi}
\subsection{Kimlik Dogrulama, Yetkilendirme ve Hesap Verebilirlik (AAA)}
AAA modeli, kimligin dogrulanmasi, yetkinin sinirlandirilmasi ve her islemin izlenebilirligini birlikte ele alir.

\subsection{Windows Active Directory ve LDAP Mimarilerinde Guvenlik Tasarimi}
AD/LDAP ortamlarinda tiering, yonetici hesap ayrimi, denetim kaydi ve MFA zorunlulugu temel tasarim prensipleridir.

\section{Formel Erisim Kontrol Modelleri}
\subsection{Istege Bagli (DAC), Zorunlu (MAC) ve Rol Tabanli (RBAC) Erisim}
RBAC operasyonel pratiklik saglarken, kritik ortamlarda MAC politikasi daha kati guvenlik ihtiyacini karsilar.

\subsection{Guvenlik Ozellikleri: Bell-LaPadula ve Biba Modelleri (Star Property Kavramlari)}
Bell-LaPadula gizlilik odakli, Biba butunluk odaklidir. Regule ortamlarda model temelli erisim tasarimi denetimlerde izlenebilirlik saglar.

\section{Ayricalikli Erisim Yonetimi (PAM)}
\subsection{Yonetici Hesaplarinin Korunmasi ve Oturum Izolasyonu}
PAM kasasi, JIT/JEA yaklasimi ve oturum kaydi bir arada kullanildiginda yanal hareket riski belirgin azalir.

\section{Modern Kimlik Dogrulama Mekanizmalari}
\subsection{Cok Faktorlu Kimlik Dogrulama (MFA), SSO ve Federasyon (SAML)}
SAML/OIDC tabanli federasyon, parola yuzeyini azaltir. MFA seciminde phishing dayanikli yontemler (FIDO2/WebAuthn) onceliklendirilmelidir.

\section{Sifir Guven (Zero Trust) Mimarisi}
\subsection{Cihaz ve Kullanici Dogrulamasina Dayali Mikro-Erisim}
Zero Trust, surekli dogrulama ve mikro erisim modeliyle "bir kez giris = tam guven" varsayimini terk eder.
\begin{figure}[htbp]
\centering
\fbox{\parbox{0.8\textwidth}{\centering Placeholder: \texttt{img/ch4-zero-trust-access-flow.png}}}
\caption{Kimlik tabanli mikro-erisim akis modeli}
\label{fig:ch4-zt-access}
\end{figure}
